% ============================================================
% Conclusion
% Synthesises the findings reported in chapters/results.tex against
% the research questions stated in the introduction. Deliberately
% avoids restating the full statistics, which belong to Section
% \ref{subsec:res-summary}; this chapter argues what they mean.
% ============================================================

\section{Conclusion}
\label{sec:conclusion}

This thesis asked whether language invariance is a useful \emph{objective} for multilingual reasoning,
rather than only a property worth measuring. The setting was multilingual numerical claim verification
as posed by CLEF~CheckThat!~2026 Task~2~\cite{clef-checkthat:2026:task2}: given a numerical claim, its
evidence, and a set of candidate reasoning traces, a system must rank the traces by usefulness and
derive a claim-level verdict, in English, Spanish, and Arabic. Starting from the supervised trace
scorer developed for our shared-task submission~\cite{khawaja2026gipplab}, we added a score-level
language-invariance objective that penalizes a model when it assigns different relevance profiles to
the same reasoning traces across languages, and an evaluation protocol that measures cross-lingual
consistency separately from task accuracy. Across $28$ non-degenerate configurations spanning three
open-weight backbones, evaluated on a held-out trilingual test set of $4{,}233$ parallel claims, the
answer is affirmative but narrower and more specific than the question invites.

\subsection{Answers to the Research Questions}
\label{subsec:concl-rqs}

\paragraph{RQ1: Does optimizing for language-invariant reasoning improve ranking performance?}
Not in the sense of retrieving better evidence, but decisively in the sense of retrieving the
\emph{same} evidence across languages. On Recall@5, the official ranking metric, invariance training
is indistinguishable from standard multilingual fine-tuning; the best Recall@5 in the study belongs to
a standard fine-tune, and both approaches sit at roughly $90\%$ of the metric's attainable ceiling.
What the objective changes is the stability of the ordering: cross-lingual rank agreement between
English and Arabic rises roughly fourfold, and the invariance-trained and non-invariance populations
do not overlap at all on this measure. A modest gain in verdict quality comes with it. The correct
claim is therefore about robustness rather than about ranking accuracy, and
Section~\ref{subsec:res-invariance} states it in those terms.

This is a less dramatic answer than ``the objective makes the ranker better,'' but it is the more
useful one. The ranking component of this task is close to saturated: with $13.7\%$ of claims having
no gold-matching trace at all, further Recall@5 gains must come from generating or retrieving better
candidate traces, not from ranking the existing ones more cleverly. Against that ceiling, an
intervention that leaves retrieval unchanged while making it reproducible across languages is
addressing the dimension that still has headroom.

\paragraph{RQ2: What is the relationship between language invariance and reasoning quality?}
Positive, moderate, and---importantly---not strong enough to treat either property as a proxy for the
other. Verdict quality correlates with cross-lingual rank agreement at $r = 0.60$, but at essentially
fixed verdict quality the observed agreement still spans more than half the range of Kendall's $\tau$.
The ranking metric is a weaker predictor still. Two observations from this thesis sharpen the point
beyond a correlation coefficient. First, the random-ranking baseline is \emph{perfectly}
language-invariant by construction and simultaneously useless, which shows that consistency alone
certifies nothing. Second, and less obviously, every system we evaluated---including the weakest---
achieves high verdict agreement across languages, while rank agreement varies by almost an order of
magnitude. Models routinely reach the same verdict in all three languages by way of largely disjoint
evidence.

That asymmetry has a methodological consequence that we consider one of the more transferable results
of this work: had we followed the common practice of measuring cross-lingual consistency only over
final answers, we would have concluded that all our systems were already language-invariant, and the
effect this thesis reports would have been invisible. Agreement statistics computed over discrete
outputs are a necessary but badly insufficient instrument, and they should be paired with
agreement over the intermediate structure a system produces---here, the trace ranking---wherever that
structure exists. This concretely extends, to a ranking setting, the accuracy/consistency decoupling
established for factual prediction~\cite{elazar2021measuringimprovingconsistencypretrained,Qi_2023}.

\paragraph{RQ3: Can language-invariant training objectives lead to more robust multilingual reasoning
models?} Yes, and the gains land where the multilingual difficulty actually is. Spanish behaves almost
identically to English for every system, so the cross-lingual problem in this task is essentially an
Arabic problem---consistent with the expectation that a different numeral script and weaker
pre-training representation would make Arabic the hard case~\cite{reddy2026effectscriptsformatsllm}.
Within Arabic, the damage is concentrated in the minority \textsf{Conflicting} class, the label that
requires genuinely weighing contradictory quantitative evidence rather than matching a single number;
this is where standard fine-tuning degrades sharply and invariance training holds. The effect appears
on both the Llama and Qwen backbones, which argues that it is a property of the objective rather than
of one model family, and it was obtained under a \emph{reduced} training-context budget, making the
comparison conservative.

The value-aware numerical components, carried over from the shared-task system, did not add to this.
Within the invariance-trained group, combining plain-text numeric normalization with the invariance
objective cost cross-lingual stability and bought no accuracy, although normalization remained mildly
positive for standard fine-tunes. The most plausible reading is redundancy rather than harm: both
mechanisms attack the same underlying problem---that the same quantity is represented differently
across scripts---and once the score-level constraint is imposed, restating canonical values in the
input adds length without adding signal. On this evidence the two should not be combined, which is a
useful negative result for the sub-thread on numerical representation.

\subsection{Contributions Revisited}
\label{subsec:concl-contrib}

Against the contributions claimed in Section~\ref{sec:contributions}, this thesis delivers the
following. It formulates language invariance as a differentiable training objective on a competitive
trace-scoring pipeline, defined at the level of per-trace scores rather than discrete outputs, and
shows that this formulation is what makes the property optimizable at all: a scorer exposes a
continuous quantity that can be aligned across languages, whereas generative rankers expose only an
ordering that must first be reconciled. It supplies a systematic comparison of ranking formulations
under a consistency lens, establishing that prompted listwise and pairwise ranking and
preference-based ranking all inherit order sensitivity that supervised trace scoring avoids. It tests
value-aware numerical representations in combination with the invariance objective and reports the
negative interaction described above. Finally, it contributes an evaluation protocol that reports task
metrics and cross-lingual agreement independently, at both the verdict and the ranking level, and
demonstrates concretely why the ranking level is the one that discriminates.

We also record two smaller findings that were not anticipated. The first is the metric-redundancy
result: across configurations, Recall@5, Precision@5, MRR@5, macro-F1, and the majority-class F1
scores are nearly collinear, so the shared task's headline numbers largely restate a single ordering
of systems, and the minority \textsf{Conflicting} class is what actually separates them. The second is
the diagnostic value of a trivially invariant baseline: including a seeded random ranker, which is
invariant by construction, provides a floor that any consistency claim must clear and makes the
independence of consistency from quality immediately legible.

\subsection{Limitations}
\label{subsec:concl-limits}

The threats to validity are stated in full in Section~\ref{subsec:res-threats}; two of them bound
the conclusions above and are worth restating plainly. The invariance-trained models were trained
under a smaller context budget than several of the standard fine-tunes, so training condition and
context budget are partially entangled. The direction of that confound is conservative---the
handicapped arm is the one that wins---and the measured effect of extra context is small and not even
consistent in sign across languages, but the comparison is observational rather than a controlled
ablation. Second, and most fundamentally, all our invariance measures are behavioural. We show that the
objective changes what a model \emph{does} across languages; we do not show that it changes how the
model represents the underlying content, and behavioural agreement is known not to entail shared
internal representation~\cite{ifergan2024beneathsurfaceconsistencyexploring}. We are also careful not
to claim a uniform win: on Spanish macro-F1 the best short-context standard fine-tune outperforms
every invariance-trained configuration, and the aggregate Arabic gap, while consistent in direction,
does not reach statistical significance on its own.

\subsection{Future Work}
\label{subsec:concl-future}

Several directions follow directly. The single most valuable next experiment is also the cheapest:
running one invariance-trained configuration at the full context budget would resolve the confound
described above, and a single seed would suffice as a sanity check. Beyond that, extending the
language set to additional scripts and typological families would test whether the effect generalizes
or is specific to the Arabic case.

A second direction concerns mechanism. Because our evidence is behavioural, pairing the objective with
representational analysis---probing or representational-similarity methods applied to the trace scorer's
hidden states~\cite{wendler2024llamasworkenglishlatent,STOLLE2026100250}---would test whether
score-level alignment actually pulls the internal representations of parallel inputs together, or
whether it merely constrains the output layer. That distinction matters for the claim that invariance
functions as an inductive bias rather than as a surface constraint.

A third direction targets the ceiling. Since ranking is nearly saturated on this candidate set, the
natural complement to this work is on the generation side: producing candidate reasoning traces with
better coverage of the gold verdict, so that a well-behaved ranker has something to find. The $13.7\%$
of claims with no gold-matching trace place a hard cap on Recall@5 that no verifier can lift.
Relatedly, the score-level invariance objective is not specific to fact-checking---it applies wherever
a model assigns comparable scores to parallel items---so testing it on other multilingual ranking and
reranking tasks would establish how far the result travels.

\subsection{Closing Remarks}
\label{subsec:concl-closing}

The clearest way to state what this thesis found is to separate two things that are easily conflated.
Standard multilingual fine-tuning already produces models that reach the same verdict in every
language; what it does not produce is models that reach it for the same reasons. Training explicitly
for language invariance does not make the underlying ranker better at finding relevant evidence, and
on the official ranking metric it changes nothing measurable. What it does is make the evidence a
model relies on transfer across languages, so that a claim and its translation are not merely assigned
the same label but are actually assessed the same way. In a fact-checking setting, where the
justification is at least as consequential as the verdict, that is the property worth having---and,
because it is invisible to the verdict-level agreement statistics ordinarily reported, it is a property
that will keep going unmeasured until the evaluation protocol is changed to look for it.
